\documentclass[10pt,openany]{book}
\usepackage{cnbook}

\renewcommand{\seriesnum}{I}
\renewcommand{\BookShortTitle}{集中不等式}
\renewcommand{\BookFooterLabel}{@MIStatlE}

\begin{document}

\MakeBookCoverUltraSimple
  {集中不等式}
  {高维概率与机器学习理论中的浓缩现象}
  {上册 · 预印本}

\frontmatter
\setcounter{tocdepth}{2}
\TOCwithoutHeadFoot

\mainmatter

\part{预备知识与工具}

\chapter{次高斯变量与 Orlicz 范数}
\ChapterEpigraph[Kolmogorov]{The theory of probability as a mathematical discipline can and should be developed from axioms in exactly the same way as geometry and algebra.}

\begin{SideBar}
\textbf{阅读建议}：这一章先建立统一语言。之后的 Bernstein、Hoeffding、Hanson--Wright 都会不断回到这些范数、矩母函数和尾部控制工具。
\end{SideBar}

\section{核心定义}

\begin{definition}[Orlicz 范数]
随机变量 $Z$ 的 $\psi_2$ 范数定义为
\[
\norm{Z}_{\psi_2} := \inf\left\{ s>0 : \E \exp\left(\frac{Z^2}{s^2}\right) \le 2 \right\}.
\]
它衡量的是尾部衰减速度，而不是单纯的方差大小。
\end{definition}

\begin{lemma}[次高斯平方的次指数性]
若 $X$ 是中心化次高斯随机变量，则 $X^2-\E X^2$ 是次指数随机变量，且其 $\psi_1$ 范数由 $\norm{X}_{\psi_2}^2$ 控制。
\end{lemma}

\begin{KeyBox}
\textbf{要点}：在高维概率中，$\psi_2$ 控制小偏差，$\psi_1$ 控制平方后的重尾项。把“次高斯 $\to$ 次指数”的链条记住，很多浓缩不等式都会自然展开。
\end{KeyBox}

\section{从矩母函数到尾界}

对独立中心化变量族 $\{X_i\}_{i=1}^n$，如果存在常数 $K$ 使得 $\norm{X_i}_{\psi_2}\le K$，那么任意线性组合都会表现出稳定的高斯型尾部。这个现象可以理解为：即使维度增大，真正主导偏差的仍然是“方向上的能量”和“单坐标上的最坏情况”。

\begin{theorem}[Bernstein 型界]
设 $Z_1,\dots,Z_n$ 独立且中心化，并满足 $\norm{Z_i}_{\psi_1}\le v$。则对任意系数向量 $a\in\mathbb{R}^n$ 与任意 $t\ge 0$，
\[
\mathbb{P}\left\{ \abs*{\sum_{i=1}^n a_i Z_i} \ge t \right\}
\le 2\exp\left[
  -c \min\left(
    \frac{t^2}{v^2\norm{a}_2^2},
    \frac{t}{v\norm{a}_\infty}
  \right)
\right].
\]
\end{theorem}

\begin{Takeaway}
小偏差受 $\norm{a}_2$ 主导，大偏差受 $\norm{a}_\infty$ 主导。这个二阶段结构几乎就是 Bernstein 不等式的直觉本体。
\end{Takeaway}

\begin{remark}
如果你在读算法分析，可以把这里的“$\norm{a}_\infty$ 控制极端坐标”理解成最坏样本或最坏梯度分量的影响。
\end{remark}

\chapter{矩阵浓缩的入口}

\section{为什么要看矩阵版本}

很多机器学习问题中的对象并不是标量和，而是样本协方差矩阵、Hessian 或随机特征矩阵。标量浓缩只告诉我们某个方向上的偏差，而矩阵浓缩试图统一控制全部方向。

\begin{definition}[样本协方差]
设 $X_1,\dots,X_n\in\mathbb{R}^d$ 独立同分布，且 $\E X_i = 0$，则样本协方差矩阵定义为
\[
\widehat{\Sigma} := \frac{1}{n}\sum_{i=1}^n X_i X_i^\top.
\]
\end{definition}

\begin{proposition}[各向同性次高斯向量]
若 $X$ 是各向同性次高斯向量，则对足够大的 $n$，$\widehat{\Sigma}$ 以高概率接近单位阵，偏差尺度大约是 $\sqrt{d/n}$ 与更高阶尾项的组合。
\end{proposition}

\begin{Example}
当 $X$ 的坐标独立且服从标准高斯时，$\widehat{\Sigma}$ 的谱范数偏差可以用矩阵 Bernstein 或更精细的随机矩阵工具来控制。这一页里的 Example 环境适合放“具体计算、反例或 sanity check”。
\end{Example}

\part{主结果}

\chapter{Hanson--Wright 不等式}
\section{二次型的浓缩}

\begin{theorem}[Hanson--Wright]
设 $X=(X_1,\dots,X_n)$ 的坐标独立、中心化且次高斯，$A\in\mathbb{R}^{n\times n}$。则存在常数 $c>0$，使得对任意 $t\ge 0$，
\[
\begin{aligned}
\mathbb{P}\left\{
  \abs*{X^\top A X - \E(X^\top A X)} \ge t
\right\}
&\le 2\exp\Biggl(
-c \min\Bigl\{
\frac{t^2}{K^4 \norm{A}_F^2}, \\
&\qquad\qquad
\frac{t}{K^2 \norm{A}_{\mathrm{op}}}
\Bigr\}
\Biggr).
\end{aligned}
\]
\end{theorem}

\begin{corollary}[谱初始化中的应用]
当随机矩阵来自噪声污染的特征模型时，Hanson--Wright 能把二次型波动压到可控范围，从而解释谱初始化和某些非凸方法为什么能稳定起步。
\end{corollary}

\begin{Takeaway}
把二次型看成“线性浓缩 + 交叉项控制”的升级版，会比死记公式更有用。真正重要的是 Frobenius 范数和算子范数分别对应的两个尺度。
\end{Takeaway}

\appendix
\part{附录}

\chapter{记号表}
\begin{itemize}
  \item $\E$：数学期望。
  \item $\Var$：方差算子。
  \item $\norm{\cdot}_{\psi_1}, \norm{\cdot}_{\psi_2}$：次指数 / 次高斯 Orlicz 范数。
  \item $\norm{A}_{\mathrm{op}}$：矩阵算子范数。
\end{itemize}

\backmatter
\begin{thebibliography}{9}
\bibitem{vershynin}
Roman Vershynin. \emph{High-Dimensional Probability}. Cambridge University Press, 2018.

\bibitem{wainwright}
Martin Wainwright. \emph{High-Dimensional Statistics}. Cambridge University Press, 2019.
\end{thebibliography}

\end{document}
